\documentclass{article}

\title{Ch. 26 CW}
\author{Principles of Macroeconomics}
\date{\today}

\begin{document}

\maketitle

\section*{Instructions}

Please do the following problems on another sheet of paper and turn them in on Gradescope before the next class.

\section*{Problems}

\begin{enumerate}

\item The accompanying table shows the average annual growth rate in real GDP per capita for Argentina, Ghana, and South Korea.  For each country, calculate how long it would take for real GDP per capita to double if their growth rates from 1990-2000 continued indefinitely.

	\begin{center}
	\begin{tabular}{| c | c | c | c |}
	\hline
	Years & Argentina & Ghana & South Korea \\
	\hline
	1990-2000 & 4.27\% & 1.33\% & 4.78\% \\
	\hline
	\end{tabular}
	\end{center}
	
\item Over the next 100 years, real GDP per capita in Groland is expected to grow at an average annual rate of 2.0\%. In Sloland, however, growth is expected to be somewhat slower, at an average annual growth rate of 1.5\%. If both countries have real GDP per capita today of \$20,000, how will their real GDP per capita differ in 100 years?

\item Which of the following is the better predictor of a future high long-run growth rate: a high standard of living today or high levels of savings and investment spending? Explain your answer. 

\item In many developing nations, young women have lower enrollment rates in secondary school than young men. Describe several ways in which greater educational opportunities for young women could lead to faster economic growth in these countries.

\item The international Property Right Index scores countries based on their legal and political environments and the extent to which they protect property rights. Go online and find a recent ranking. Choose three countries with high scores and three countries with low scores. Then find estimates of GDP per person in each of these six countries. What pattern do you find? Give two possible interpretations of the pattern.

\item International data show a positive correlation between income per person and the health of the population.

	\begin{enumerate}
	
	\item Explain how higher income might cause better health outcomes.
	
	\item Explain how better health outcomes might cause higher income.
	
	\item How might the relative importance of your two hypotheses be relevant for public policy?
	
	\end{enumerate}
	
\item The great 18th-century economist Adam Smith wrote, "Little else is requisite to carry a state to the highest degree of opulence from the lowest barbarism but peace, easy taxes, and a tolerable administration of justice: all the rest being brought about by the natural course of things." Explain how each of the three conditions Smith describes promotes economic growth. 

\end{enumerate}

\end{document}